\section{An invariant-kernel obstruction}
\label{sec:obstruction}

We first prove the criterion announced in the introduction. A large
Hilbert--Schmidt energy will force an operator far from the translation
commutant; a sequence of such operators with uniformly bounded
conjugation differences will then give one bounded non-inner cocycle.
Write $\mathcal B(H)$ for the bounded operators on $H$ and
$\norm{X}_{\HS}^2=\Tr(X^*X)$ for the Hilbert--Schmidt norm of a
finite matrix, with the unnormalized trace $\Tr$.
All operator norms below are ordinary Hilbert-space operator norms.
The finite-dimensional fibres will enter through
Hilbert--Schmidt estimates for one block row and one block column.

For a countable group $G$ and a positive integer $k$, put
\[
  \mathcal H_k=\ell^2(G;\C^k),
  \qquad (U_k(g)\xi)(x)=\xi(g^{-1}x).
\]
For $x\in G$, let $\iota_x:\C^k\to\mathcal H_k$ be the isometric
inclusion into the coordinate indexed by $x$.
For $F\in\mathcal B(\mathcal H_k)$, write
\[
  F[x,y]=\iota_x^*F\iota_y,\qquad
  \operatorname{row}_e(F)=\iota_e^*F,\qquad
  \operatorname{col}_e(F)=F\iota_e.
\]
Thus a block row has codomain $\C^k$ and a block column has domain
$\C^k$. Write $U_k(G)'$ for the bounded operators commuting with
every $U_k(g)$. If $F\in U_k(G)'$, then
$F[x,y]=F[g^{-1}x,g^{-1}y]$, because $U_k(g)\iota_y=\iota_{gy}$.
Taking $g=x$ gives
\[
  F[x,y]=V(x^{-1}y),\qquad V(s)=F[e,s].
\]
This statement concerns the blocks of a bounded operator and requires
no summability assumption on a convolution formula.

\Needspace{18\baselineskip}
\begin{lemma}[Distance from the translation commutant]
\label{lem:kernel-obstruction}
Let $G$ be a countable group, let $k$ be a positive integer, and let
$A,T\in\mathcal B(\mathcal H_k)$ with $A\in U_k(G)'$.
Write
\[
  A[x,y]=W(x^{-1}y),\qquad
  E=\sum_{s\in G}\norm{W(s)}_{\HS}^2.
\]
Suppose that, for some $K\ge0$,
\[
  \norm{\operatorname{row}_e(T)}\le K,
  \qquad
  \norm{\operatorname{col}_e(A-T)}\le K.
\]
Then $E$ is finite and, for every $C\in U_k(G)'$,
\begin{equation}
\label{eq:kernel-energy-obstruction}
  E\le 4k\bigl(K+\norm{T-C}\bigr)^2.
\end{equation}
Consequently,
\begin{equation}
\label{eq:kernel-distance-obstruction}
  \dist\bigl(T,U_k(G)'\bigr)
  \ge \frac12\sqrt{\frac{E}{k}}-K,
\end{equation}
where $\dist(T,U_k(G)')=\inf_{C\in U_k(G)'}\norm{T-C}$.
\end{lemma}

\begin{proof}
Fix $C\in U_k(G)'$, and put
\[
  B=T-C,\qquad h=K+\norm{B},\qquad Q(s)=C[e,s].
\]
The kernel identity above gives $C[x,y]=Q(x^{-1}y)$.
Since $C=T-B$, the operator
\[
  R=\operatorname{row}_e(C):\mathcal H_k\longrightarrow\C^k
\]
has norm at most $h$.
Any bounded operator with codomain $\C^k$ is Hilbert--Schmidt and
has squared Hilbert--Schmidt norm at most $k$ times its squared
operator norm: indeed,
\[
  \norm{R}_{\HS}^2=\Tr(RR^*)\le k\norm{R}^2.
\]
Computing this norm using the coordinate orthonormal basis of
$\mathcal H_k$ therefore gives
\begin{equation}
\label{eq:kernel-row-energy}
  \sum_{s\in G}\norm{Q(s)}_{\HS}^2
  =\norm{R}_{\HS}^2\le kh^2.
\end{equation}

On the other hand, $A-C=(A-T)+B$, so
\[
  L=\operatorname{col}_e(A-C):\C^k\longrightarrow\mathcal H_k
\]
also has norm at most $h$.
Its domain has dimension $k$, whence
\[
  \norm{L}_{\HS}^2=\Tr(L^*L)\le k\norm{L}^2\le kh^2.
\]
The block of this column at row $x$ is
\[
  (A-C)[x,e]=W(x^{-1})-Q(x^{-1}).
\]
Inversion is a bijection of $G$; computing the Hilbert--Schmidt norm
by blocks and reindexing therefore gives
\begin{equation}
\label{eq:kernel-column-energy}
  \sum_{s\in G}\norm{W(s)-Q(s)}_{\HS}^2
  =\norm{L}_{\HS}^2\le kh^2.
\end{equation}
In both computations, the sum of the squared block norms is exactly
the sum of squared matrix coefficients defining the
Hilbert--Schmidt norm.  The group is countable, and all these
summands are nonnegative.

By \Cref{eq:kernel-row-energy,eq:kernel-column-energy}, both $Q$
and $W-Q$ are square summable in the Hilbert--Schmidt norm.
The triangle inequality for these coefficient families gives
\[
  \left(\sum_{s\in G}\norm{W(s)}_{\HS}^2\right)^{1/2}
  \le \sqrt{k}\,h+\sqrt{k}\,h=2\sqrt{k}\,h.
\]
This proves finiteness of $E$ and
\Cref{eq:kernel-energy-obstruction}.
Rearranging that inequality and then taking the infimum over $C$
proves \Cref{eq:kernel-distance-obstruction}.
\end{proof}

The next proposition turns an unbounded sequence of these distances
into one uniformly bounded representation that cannot be
unitarized.  The relation between bounded operator cocycles and
triangular representations is standard; see
\cite[Proposition~2.4]{EpsteinMonod2009},
\cite[Section~2]{MonodOzawa2010}, and
\cite[Remark~4 and Proposition~5(ii)]{Pisier2011}.
We include the argument, in particular the bounded-complement step
and the direct-sum quantifiers.

\begin{proposition}[Direct-sum obstruction]
\label{prop:direct-sum-obstruction}
Let $G$ be a countable group.
For each $j\in\N$, let $k_j$ be a positive integer, put
\[
  \mathcal H_j=\ell^2(G;\C^{k_j}),
  \qquad (U_j(g)\xi)(x)=\xi(g^{-1}x),
\]
and let $A_j,T_j\in\mathcal B(\mathcal H_j)$ with
$A_j\in U_j(G)'$.
Write
\[
  A_j[x,y]=W_j(x^{-1}y),\qquad
  E_j=\sum_{s\in G}\norm{W_j(s)}_{\HS}^2.
\]
Suppose that one constant $K\ge0$ satisfies
\begin{align*}
  \norm{\operatorname{row}_e(T_j)}&\le K, &
  \norm{\operatorname{col}_e(A_j-T_j)}&\le K
  &&(j\in\N),\\
  \norm{T_j-U_j(g)T_jU_j(g)^{-1}}&\le K
  &&&&(j\in\N,\ g\in G).
\end{align*}
Assume also that
\[
  \frac{E_j}{k_j}\longrightarrow\infty.
\]
Then $G$ has a representation on a separable complex Hilbert space
whose operators are bounded and invertible, with
\[
  \sup_{g\in G}\norm{\pi(g)}\le1+K,
\]
and which is not similar to a unitary representation.
No uniform bound on $\norm{T_j}$ or $\norm{A_j}$ is assumed.
\end{proposition}

\begin{proof}
On the Hilbert direct sum
\[
  \mathcal H=\bigoplus_{j\in\N}\mathcal H_j
\]
define
\[
  U(g)=\bigoplus_{j\in\N}U_j(g),\qquad
  D_j(g)=T_j-U_j(g)T_jU_j(g)^{-1},\qquad
  D(g)=\bigoplus_{j\in\N}D_j(g).
\]
The operator $U(g)$ is unitary and $\norm{D(g)}\le K$ by the
hypotheses.  Thus these direct sums are bounded operators on all of
$\mathcal H$.  No bounded operator $\bigoplus_jT_j$ is needed.
Each $\mathcal H_j$ is separable because $G$ is countable and $k_j$
is finite; consequently $\mathcal H$ and $\mathcal H\oplus\mathcal H$
are separable.

Expanding the definition of $D_j$ gives
\[
  D_j(gh)=D_j(g)+U_j(g)D_j(h)U_j(g)^{-1}.
\]
Taking bounded direct sums gives the cocycle identity
\begin{equation}
\label{eq:direct-sum-cocycle}
  D(gh)=D(g)+U(g)D(h)U(g)^{-1}
  \qquad(g,h\in G).
\end{equation}
Define an operator on $\mathcal H\oplus\mathcal H$ by
\begin{equation}
\label{eq:triangular-representation}
  \pi(g)=
  \begin{pmatrix}
    U(g)&D(g)U(g)\\
    0&U(g)
  \end{pmatrix}.
\end{equation}
By \Cref{eq:direct-sum-cocycle}, its upper-right block satisfies
\[
  D(gh)U(gh)
  =D(g)U(g)U(h)+U(g)D(h)U(h).
\]
Multiplication of block matrices now shows
$\pi(g)\pi(h)=\pi(gh)$.
Also $D(e)=0$, so $\pi(e)=\one$.
It follows that each $\pi(g)$ is invertible, with
$\pi(g)^{-1}=\pi(g^{-1})$.
The factorization
\[
  \pi(g)=
  \begin{pmatrix}\one&D(g)\\0&\one\end{pmatrix}
  \begin{pmatrix}U(g)&0\\0&U(g)\end{pmatrix}
\]
has a unitary second factor, whereas the first factor is the
identity plus an operator of norm $\norm{D(g)}$.
Therefore $\norm{\pi(g)}\le1+\norm{D(g)}\le1+K$ for every $g$.

Suppose, for a contradiction, that a bounded invertible operator
$S$ makes $V(g)=S\pi(g)S^{-1}$ unitary for all $g\in G$.
The closed subspace $M=\mathcal H\oplus0$ is invariant under every
$\pi(g)$ and its inverse.  Consequently $SM$ is closed and
invariant under every $V(g)$ and its inverse, so $(SM)^\perp$ is
also invariant.  The subspace
\[
  N=S^{-1}\bigl((SM)^\perp\bigr)
\]
is thus a closed invariant complement of $M$:
$\mathcal H\oplus\mathcal H=M\mathbin{\dotplus}N$.

Let $P_1,P_2:\mathcal H\oplus\mathcal H\to\mathcal H$ be the
coordinate projections.
The restriction $P_2|_N:N\to\mathcal H$ is injective, because
$N\cap M=0$, and surjective, because every $(0,\eta)$ has a
decomposition into an element of $M$ and an element of $N$.
It is a bounded bijection between Banach spaces, hence has bounded
inverse.  Therefore
\[
  B=P_1(P_2|_N)^{-1}\in\mathcal B(\mathcal H),
  \qquad
  N=\{(B\eta,\eta):\eta\in\mathcal H\}.
\]
Applying $\pi(g)$ to this graph and using its invariance gives
\[
  U(g)B\eta+D(g)U(g)\eta=BU(g)\eta
  \qquad(\eta\in\mathcal H).
\]
Since $U(g)$ is invertible, this is equivalent to
\begin{equation}
\label{eq:global-cocycle-implementer}
  D(g)=B-U(g)BU(g)^{-1}.
\end{equation}

Let $I_j:\mathcal H_j\to\mathcal H$ be the canonical isometric
inclusion, and set $B_j=I_j^*BI_j$.
The summand $\mathcal H_j$ reduces every $U(g)$, so
\[
  I_j^*U(g)BU(g)^{-1}I_j
  =U_j(g)B_jU_j(g)^{-1}.
\]
Compressing \Cref{eq:global-cocycle-implementer} therefore yields
\[
  D_j(g)=B_j-U_j(g)B_jU_j(g)^{-1},
  \qquad\norm{B_j}\le\norm{B}.
\]
This compression does not require $B$ to preserve the summands.
Comparison with the definition of $D_j$ shows that
\[
  C_j=T_j-B_j\in U_j(G)'.
\]
Apply \Cref{lem:kernel-obstruction} to $A_j,T_j$ and this choice of
$C_j$.  Its energy estimate gives, for every $j$,
\[
  E_j\le4k_j\bigl(K+\norm{B_j}\bigr)^2
      \le4k_j\bigl(K+\norm{B}\bigr)^2.
\]
This contradicts $E_j/k_j\to\infty$, because $B$ is one bounded
operator.  Hence the representation in
\Cref{eq:triangular-representation} is not similar to a unitary
representation.
\end{proof}

For clarity, each individual $D_j$ in this construction is already
implemented by the bounded operator $T_j$.
The obstruction is to a bounded implementer on the whole direct
sum.  Indeed, \Cref{lem:kernel-obstruction} shows that any operator
$B_j$ satisfying
$D_j(g)=B_j-U_j(g)B_jU_j(g)^{-1}$ for all $g$ must obey
\[
  \norm{B_j}\ge\frac12\sqrt{\frac{E_j}{k_j}}-K.
\]
Thus the proof requires only the existence of a similarity for one
representation in order to obtain a contradiction; it assumes no
uniform bound on similarity operators across different
representations.
